\providecommand{\C}{\mathbb{C}}
\providecommand{\Z}{\mathbb{Z}}
\providecommand{\cO}{\mathcal{O}}
\providecommand{\fg}{\mathfrak{g}}
\providecommand{\At}{\mathrm{At}}
\providecommand{\End}{\mathrm{End}}
\providecommand{\Sym}{\mathrm{Sym}}
\providecommand{\Tr}{\mathrm{Tr}}
\providecommand{\PV}{\mathrm{PV}}
\providecommand{\BV}{\mathrm{BV}}
\providecommand{\BCOV}{\mathrm{BCOV}}
\providecommand{\hCS}{\mathrm{hCS}}
\providecommand{\kcat}{\kappa_{\mathrm{cat}}}
\providecommand{\kchHodge}{\kappa_{\mathrm{ch}}^{\mathrm{Hodge}}}
\providecommand{\kchHeis}{\kappa_{\mathrm{ch}}^{\mathrm{Heis}}}
\providecommand{\kchBCOV}{\kappa_{\mathrm{ch}}^{\mathrm{BCOV}}}
\providecommand{\kBKM}{\kappa_{\mathrm{BKM}}}
\providecommand{\kfiber}{\kappa_{\mathrm{fiber}}}

\section*{Atiyah class, Yukawa cubic, BCOV anomaly, and holomorphic
Chern-Simons on compact CY$_3$}
\label{sec:bcov-yukawa-hcs-cy3}

Let $X$ be a compact Calabi-Yau threefold with holomorphic volume form
$\Omega_X$ and let $\fg$ be a finite-dimensional Lie algebra equipped
with an invariant non-degenerate symmetric pairing. Four structures enter
the B-model and holomorphic Chern-Simons calculus:
\[
\At(T_X),\qquad
\ell_\bullet^{\mathrm{Kap}},\qquad
Y_3,\qquad
\alpha_{\BCOV}.
\]
They have different cohomological homes. The Atiyah class is a vector
bundle extension class. Kapranov brackets are the induced
$L_\infty$-operations on shifted tangent polyvectors. The Yukawa cubic is
the genus-zero three-point tensor on complex-structure deformations. The
BCOV curving is the one-loop Atiyah-dilaton anomaly of the BV measure.
The holomorphic Chern-Simons BV theory uses all four data after a choice
of gauge fixing, orientation, renormalisation scheme, factorisation
algebra, and completion.

\subsection*{The Atiyah class}
\label{subsec:bcov-yukawa-atiyah}

\begin{definition}[Atiyah class]
\label{def:bcov-yukawa-atiyah}
For a holomorphic vector bundle $E$ on a complex manifold $X$, the
Atiyah class is the extension class of the first jet sequence
\[
0\longrightarrow \Omega^1_X\otimes E
\longrightarrow J^1(E)
\longrightarrow E
\longrightarrow 0.
\]
For the tangent bundle it is
\[
\At(T_X)\in
\operatorname{Ext}^1_{\cO_X}(T_X,\Omega^1_X\otimes T_X)
=H^1(X,\Omega^1_X\otimes \End T_X).
\]
It obstructs a holomorphic connection on $T_X$. Its trace is
$\Tr\At(T_X)=\At(\det T_X)$, hence the trace vanishes after a fixed
Calabi-Yau trivialisation of $\omega_X=\det \Omega^1_X$. The traceless
class may still be non-zero.
\end{definition}

For $X=K3\times E$ the tangent bundle splits under the two projections
$p_1,p_2$:
\[
T_X\simeq p_1^*T_{K3}\oplus p_2^*T_E,\qquad
\At(T_X)=p_1^*\At(T_{K3})\oplus p_2^*\At(T_E).
\]
The elliptic summand has $\At(T_E)=0$ because $T_E\simeq \cO_E$ has the
translation-invariant holomorphic connection. The K3 summand has
non-zero Atiyah square; its trace square integrates to $c_2(K3)=24$.
Thus $K3\times E$ is trace-trivial as a Calabi-Yau threefold while still
carrying non-flat tangent Atiyah data from the K3 factor.

\subsection*{Kapranov brackets}
\label{subsec:bcov-yukawa-kapranov}

Kapranov's construction assigns to a complex manifold a Lie algebra
object $T_X[-1]$ in the derived category. On a K\"ahler representative it
is realised by an $L_\infty$-algebra on
\[
\Omega^{0,\bullet}(X,T_X[-1]).
\]
The binary operation is governed by $\At(T_X)$; higher brackets are
obtained from covariant derivatives and iterated contractions of the
same curvature class. The $L_\infty$ identities are the Bianchi identities
for the curvature representative.

This is the precise formality statement used here. The Hochschild-Kostant-
Rosenberg map identifies Hochschild cohomology with polyvectors as a
graded vector space. On a curved compact complex manifold it requires
Atiyah and Todd corrections to respect the Gerstenhaber structure. The
strict HKR map is therefore replaced by an $L_\infty$ quasi-isomorphism
whose first non-trivial corrections are controlled by the Atiyah class.
This is the content of Kapranov's theorem and the
Calaque-Caldararu-Tu comparison of Atiyah classes with HKR.

\subsection*{The Yukawa cubic}
\label{subsec:bcov-yukawa-y3}

The Yukawa coupling of the B-model is the genus-zero cubic tensor
\[
Y_3\in \Sym^3 H^1(X,T_X)^\vee .
\]
For $v_1,v_2,v_3\in H^1(X,T_X)$ it may be written, after fixing
$\Omega_X$, as
\[
Y_3(v_1,v_2,v_3)
=
\int_X
\Omega_X\wedge
\bigl((v_1\wedge v_2\wedge v_3)\,\lrcorner\,\Omega_X\bigr).
\]
Equivalently, it is the third covariant derivative of the genus-zero
prepotential in special geometry. Under the Calabi-Yau contraction
$\wedge^3T_X\simeq\cO_X$ it lands in
\[
H^3(X,\wedge^3T_X)^\vee\simeq H^3(X,\cO_X)^\vee
\simeq H^{0,3}(X)^\vee .
\]
In the transferred Kodaira-Spencer or Kapranov minimal model, this tensor
is the tree-level cubic operation. It is a genus-zero structure constant.

For a smooth quintic threefold $X_5\subset\mathbb P^4$ with hyperplane
class $H$,
\[
Y_{HHH}^{\mathrm{cl}}(X_5)=\int_{X_5}H^3=5.
\]
This classical value is the large-volume A-model normalisation mirrored
by the B-model Yukawa tensor. It is independent of the one-loop BCOV
Euler coefficient.

\subsection*{BCOV one-loop curving}
\label{subsec:bcov-yukawa-alpha}

The Costello-Li BCOV theory is a BV theory on polyvector fields
\[
\PV^{p,q}(X)=\Omega^{0,q}(X,\wedge^pT_X).
\]
Its odd symplectic pairing is obtained by Calabi-Yau integration against
$\Omega_X$. The one-loop curving is the Atiyah-dilaton class
\[
\alpha_{\BCOV}(X)
=
\frac{\chi_{\mathrm{top}}(X)}{24}\,\delta_X
\in H^1(X,\cO_X),
\]
where $\delta_X$ denotes the structure-sheaf class obtained from the
trace of the tangent Atiyah class by the Calabi-Yau residue contraction.
The rational scalar $\chi_{\mathrm{top}}(X)/24$ is the BCOV genus-one
Euler coefficient. The cohomology class also depends on the target
$H^1(X,\cO_X)$.

Thus strict compact Calabi-Yau threefolds with $H^1(X,\cO_X)=0$ have
vanishing leading $\alpha_{\BCOV}$ as a structure-sheaf class. The
quintic still has the non-zero genus-one scalar
\[
\chi_{\mathrm{top}}(X_5)=-200,\qquad
\kchBCOV(X_5)=\frac{\chi_{\mathrm{top}}(X_5)}{24}=-\frac{25}{3}.
\]
This scalar belongs to the BCOV genus-one expansion. It is separate from
the Hodge supertrace
\[
\kchHodge(X_5)
=\sum_{q=0}^{3}(-1)^q h^{0,q}(X_5)
=1-0+0-1=0.
\]

For $X=K3\times E$,
\[
\chi_{\mathrm{top}}(K3\times E)=\chi_{\mathrm{top}}(K3)\chi_{\mathrm{top}}(E)
=24\cdot 0=0.
\]
Although $H^1(K3\times E,\cO)\simeq H^1(E,\cO_E)\simeq\C$, the leading
Atiyah-dilaton class satisfies
\[
\alpha_{\BCOV}(K3\times E)=0.
\]
The tree-level Yukawa line remains one-dimensional:
\[
H^{0,3}(K3\times E)
\simeq H^{0,2}(K3)\otimes H^{0,1}(E)
\simeq \C.
\]
The vanishing of $\alpha_{\BCOV}$ is a statement about the one-loop
Euler prefactor. It leaves the tree-level cubic tensor untouched.

\subsection*{Holomorphic Chern-Simons BV data}
\label{subsec:bcov-yukawa-hcs}

The classical holomorphic Chern-Simons field is
\[
\mathcal A\in\Omega^{0,\bullet}(X,\fg)[1],
\]
and the classical action is
\[
S_{\hCS}(\mathcal A)
=
\int_X\Omega_X\wedge
\left(
\frac{1}{2}\langle\mathcal A,\bar\partial\mathcal A\rangle
+\frac{1}{6}\langle\mathcal A,[\mathcal A,\mathcal A]\rangle
\right).
\]
The invariant pairing on $\fg$, the Jacobi identity, and the
Calabi-Yau integration pairing imply the classical master equation.
The quantum statement is stronger. A compact theorem of quantisation
must supply:
\begin{enumerate}
\item a gauge fixing and heat-kernel or equivalent propagator;
\item a determinant-line orientation and BV Laplacian compatible with
the Calabi-Yau trace;
\item renormalised counterterms solving the quantum master equation;
\item cancellation of local gauge and gravitational anomalies;
\item Costello-Gwilliam factorisation observables;
\item TCFT or BCOV background data for the open-closed coupling;
\item analytic completion of operator products and filtrations;
\item the CY$_3$ specialisation datum producing a chiral object on a
chosen curve.
\end{enumerate}

For non-abelian $\fg$ in complex dimension three, the cohomological local
gauge anomaly is governed by the quartic invariant-polynomial descent
slot, schematically
\[
I_4^\rho(x)=\Tr_\rho(x^4)
\quad\text{in}\quad
\Sym^4(\fg^\vee)^\fg .
\]
The quadratic Casimir controls the two-point function and field-strength
renormalisation. It does not supply the quartic descent primitive required
by the quantum master equation. A scalar such as
$\chi_{\mathrm{top}}(X)/24$ also does not cancel a non-abelian local
gauge anomaly without a counterterm in the same BV local functional
complex.

\subsection*{Three cohomological positions}
\label{subsec:bcov-yukawa-three-slots}

The compact CY$_3$ calculus separates the three Atiyah-sourced positions:
\[
\At(T_X)\in H^1(X,\Omega^1_X\otimes\End T_X),
\]
\[
Y_3\in\Sym^3H^1(X,T_X)^\vee
\simeq H^{0,3}(X)^\vee ,
\]
\[
\alpha_{\BCOV}(X)\in H^1(X,\cO_X)=H^{0,1}(X).
\]
The first class is the tangent-bundle extension class. The second is the
tree-level cubic tensor. The third is the one-loop BV curving. A theorem
relating them must name the map between these three receptacles. The
standard BCOV and hCS statements give a common Atiyah source and separate
genus orders; they do not identify the three classes.

\begin{proposition}[Quintic constants]
\label{prop:bcov-yukawa-quintic-constants}
For a smooth quintic threefold $X_5\subset\mathbb P^4$,
\[
\int_{X_5}H^3=5,\qquad
c(T_{X_5})=\frac{(1+H)^5}{1+5H}=1+10H^2-40H^3,
\]
\[
\chi_{\mathrm{top}}(X_5)=\int_{X_5}c_3(T_{X_5})=-200,
\]
\[
\kcat(X_5)=\chi(\cO_{X_5})=0,\qquad
\kchHodge(X_5)=0,\qquad
\kchBCOV(X_5)=-\frac{25}{3}.
\]
The number $5$ is the large-volume classical Yukawa coefficient. The
number $-25/3$ is the BCOV genus-one Euler coefficient.
\end{proposition}

\begin{proof}
Adjunction gives
$c(T_{X_5})=c(T_{\mathbb P^4}|_{X_5})/c(\cO_{X_5}(5))=(1+H)^5/(1+5H)$.
The degree-three term is $-40H^3$, and $\int_{X_5}H^3=5$. The Hodge
column is $(h^{0,0},h^{0,1},h^{0,2},h^{0,3})=(1,0,0,1)$, giving the
displayed Hodge supertrace and holomorphic Euler characteristic. The
BCOV scalar is the standard $\chi_{\mathrm{top}}/24$ genus-one Euler
coefficient.
\end{proof}

\begin{proposition}[Product constants for $K3\times E$]
\label{prop:bcov-yukawa-k3e-constants}
For $X=K3\times E$,
\[
\kcat(X)=0,\qquad
\sum_q(-1)^qh^{0,q}(X)=0,
\]
and the Vol III specialisation values are
\[
\{\kcat,\kchHeis,\kBKM,\kfiber\}(K3\times E)=\{0,3,5,24\}.
\]
Here $\kcat=0$ is the total-space holomorphic Euler characteristic,
$\kchHeis=3$ is the Heisenberg-Mukai specialisation, $\kBKM=5$ is the
Gritsenko $\Delta_5$ Borcherds weight, and $\kfiber=24$ is the K3
Mukai-lattice fibre rank. The compact Hodge supertrace of the product is
zero.
\end{proposition}

\begin{proof}
K\"unneth gives
$\chi(\cO_{K3\times E})=\chi(\cO_{K3})\chi(\cO_E)=2\cdot0=0$ and
\[
\sum_q(-1)^qh^{0,q}(K3\times E)
=(1-0+1)(1-1)=0.
\]
The values $3$, $5$, and $24$ come from three additional constructions:
the Heisenberg-Mukai chiral specialisation, the Borcherds product for
$\Delta_5$, and the K3 Mukai lattice.
\end{proof}

\subsection*{Position of the CY-to-chiral functor}
\label{subsec:bcov-yukawa-phi3}

The specialised CY-to-chiral functor in dimension $3$ produces an
$E_1$-chiral object after the choice of specialisation data
$(\Sigma_2,C)$ and completion. Any $E_2$-structure belongs to a centre,
Drinfeld double, or module-category layer. Thus an hCS or BCOV statement
on a compact CY$_3$ does not by itself construct an $E_2$ chiral algebra
on the output curve. It gives input data for an $E_1$ specialisation once
the compact BV and factorisation hypotheses above have been met.

\subsection*{Sources}
\label{subsec:bcov-yukawa-sources}

\begin{itemize}
\item Atiyah 1957, \emph{Complex analytic connections in fibre bundles}.
\item Kapranov 1999, \emph{Rozansky-Witten invariants via Atiyah classes}.
\item Calaque-Caldararu-Tu 2014, comparison of HKR and Atiyah classes.
\item Bershadsky-Cecotti-Ooguri-Vafa 1993 and 1994, holomorphic anomaly
and Yukawa couplings in the B-model.
\item Costello 2007 and Costello-Li 2012/2016, BCOV theory and BV
quantisation on Calabi-Yau manifolds.
\item Costello-Gwilliam, factorisation algebras and BV quantisation.
\end{itemize}
